%File: anonymous-submission-latex-2026.tex
\documentclass[letterpaper]{article} % DO NOT CHANGE THIS
\usepackage[submission]{aaai2026}  % DO NOT CHANGE THIS
\usepackage{times}  % DO NOT CHANGE THIS
\usepackage{helvet}  % DO NOT CHANGE THIS
\usepackage{courier}  % DO NOT CHANGE THIS
\usepackage[hyphens]{url}  % DO NOT CHANGE THIS
\usepackage{graphicx} % DO NOT CHANGE THIS
\urlstyle{rm} % DO NOT CHANGE THIS
\def\UrlFont{\rm}  % DO NOT CHANGE THIS
\usepackage{natbib}  % DO NOT CHANGE THIS AND DO NOT ADD ANY OPTIONS TO IT
\usepackage{caption} % DO NOT CHANGE THIS AND DO NOT ADD ANY OPTIONS TO IT
\frenchspacing  % DO NOT CHANGE THIS
\setlength{\pdfpagewidth}{8.5in} % DO NOT CHANGE THIS
\setlength{\pdfpageheight}{11in} % DO NOT CHANGE THIS
%
% These are recommended to typeset algorithms but not required. See the subsubsection on algorithms. Remove them if you don't have algorithms in your paper.
\usepackage{algorithm}

% Add more
\usepackage{subcaption}
\usepackage{algpseudocode} 
\usepackage{amsmath}
\usepackage{multirow}
\usepackage{amsfonts} % For \mathbb
%
% These are recommended to typeset listings but not required. See the subsubsection on listing. Remove this block if you don't have listings in your paper.
\usepackage{newfloat}
\usepackage{listings}
\DeclareCaptionStyle{ruled}{labelfont=normalfont,labelsep=colon,strut=off} % DO NOT CHANGE THIS
\lstset{%
	basicstyle={\footnotesize\ttfamily},% footnotesize acceptable for monospace
	numbers=left,numberstyle=\footnotesize,xleftmargin=2em,% show line numbers, remove this entire line if you don't want the numbers.
	aboveskip=0pt,belowskip=0pt,%
	showstringspaces=false,tabsize=2,breaklines=true}
\floatstyle{ruled}
\newfloat{listing}{tb}{lst}{}
\floatname{listing}{Listing}
%
% Keep the \pdfinfo as shown here. There's no need
% for you to add the /Title and /Author tags.
\pdfinfo{
/TemplateVersion (2026.1)
}

% DISALLOWED PACKAGES
% \usepackage{authblk} -- This package is specifically forbidden
% \usepackage{balance} -- This package is specifically forbidden
% \usepackage{color (if used in text)
% \usepackage{CJK} -- This package is specifically forbidden
% \usepackage{float} -- This package is specifically forbidden
% \usepackage{flushend} -- This package is specifically forbidden
% \usepackage{fontenc} -- This package is specifically forbidden
% \usepackage{fullpage} -- This package is specifically forbidden
% \usepackage{geometry} -- This package is specifically forbidden
% \usepackage{grffile} -- This package is specifically forbidden
% \usepackage{hyperref} -- This package is specifically forbidden
% \usepackage{navigator} -- This package is specifically forbidden
% (or any other package that embeds links such as navigator or hyperref)
% \indentfirst} -- This package is specifically forbidden
% \layout} -- This package is specifically forbidden
% \multicol} -- This package is specifically forbidden
% \nameref} -- This package is specifically forbidden
% \usepackage{savetrees} -- This package is specifically forbidden
% \usepackage{setspace} -- This package is specifically forbidden
% \usepackage{stfloats} -- This package is specifically forbidden
% \usepackage{tabu} -- This package is specifically forbidden
% \usepackage{titlesec} -- This package is specifically forbidden
% \usepackage{tocbibind} -- This package is specifically forbidden
% \usepackage{ulem} -- This package is specifically forbidden
% \usepackage{wrapfig} -- This package is specifically forbidden
% DISALLOWED COMMANDS
% \nocopyright -- Your paper will not be published if you use this command
% \addtolength -- This command may not be used
% \balance -- This command may not be used
% \baselinestretch -- Your paper will not be published if you use this command
% \clearpage -- No page breaks of any kind may be used for the final version of your paper
% \columnsep -- This command may not be used
% \newpage -- No page breaks of any kind may be used for the final version of your paper
% \pagebreak -- No page breaks of any kind may be used for the final version of your paper
% \pagestyle -- This command may not be used
% \tiny -- This is not an acceptable font size.
% \vspace{- -- No negative value may be used in proximity of a caption, figure, table, section, subsection, subsubsection, or reference
% \vskip{- -- No negative value may be used to alter spacing above or below a caption, figure, table, section, subsection, subsubsection, or reference

\setcounter{secnumdepth}{0} %May be changed to 1 or 2 if section numbers are desired.

% The file aaai2026.sty is the style file for AAAI Press
% proceedings, working notes, and technical reports.
%

% Title

% Your title must be in mixed case, not sentence case.
% That means all verbs (including short verbs like be, is, using,and go),
% nouns, adverbs, adjectives should be capitalized, including both words in hyphenated terms, while
% articles, conjunctions, and prepositions are lower case unless they
% directly follow a colon or long dash
\title{SHeRL-FL: When Representation Learning Meets Split Learning in Hierarchical Federated Learning}
\author{
    %Authors
    % All authors must be in the same font size and format.
    Written by AAAI Press Staff\textsuperscript{\rm 1}\thanks{With help from the AAAI Publications Committee.}\\
    AAAI Style Contributions by Pater Patel Schneider,
    Sunil Issar,\\
    J. Scott Penberthy,
    George Ferguson,
    Hans Guesgen,
    Francisco Cruz\equalcontrib,
    Marc Pujol-Gonzalez\equalcontrib
}
\affiliations{
    %Afiliations
    \textsuperscript{\rm 1}Association for the Advancement of Artificial Intelligence\\
    % If you have multiple authors and multiple affiliations
    % use superscripts in text and roman font to identify them.
    % For example,

    % Sunil Issar\textsuperscript{\rm 2},
    % J. Scott Penberthy\textsuperscript{\rm 3},
    % George Ferguson\textsuperscript{\rm 4},
    % Hans Guesgen\textsuperscript{\rm 5}
    % Note that the comma should be placed after the superscript

    1101 Pennsylvania Ave, NW Suite 300\\
    Washington, DC 20004 USA\\
    % email address must be in roman text type, not monospace or sans serif
    proceedings-questions@aaai.org
%
% See more examples next
}

%Example, Single Author, ->> remove \iffalse,\fi and place them surrounding AAAI title to use it
\iffalse
\title{My Publication Title --- Single Author}
\author {
    Author Name
}
\affiliations{
    Affiliation\\
    Affiliation Line 2\\
    name@example.com
}
\fi

\iffalse
%Example, Multiple Authors, ->> remove \iffalse,\fi and place them surrounding AAAI title to use it
\title{My Publication Title --- Multiple Authors}
\author {
    % Authors
    First Author Name\textsuperscript{\rm 1},
    Second Author Name\textsuperscript{\rm 2},
    Third Author Name\textsuperscript{\rm 1}
}
\affiliations {
    % Affiliations
    \textsuperscript{\rm 1}Affiliation 1\\
    \textsuperscript{\rm 2}Affiliation 2\\
    firstAuthor@affiliation1.com, secondAuthor@affilation2.com, thirdAuthor@affiliation1.com
}
\fi


% REMOVE THIS: bibentry
% This is only needed to show inline citations in the guidelines document. You should not need it and can safely delete it.
\usepackage{bibentry}
% END REMOVE bibentry

\begin{document}

\maketitle

\begin{abstract}
    Federated learning (FL) is a promising approach for addressing scalability and latency issues in large-scale networks by enabling collaborative model training without requiring the sharing of raw data. However, existing FL frameworks often overlook the computational heterogeneity of edge clients and the growing training burden on resource-limited devices. However, FL suffers from high communication costs and complex model aggregation, especially with large models. Previous works combine split learning (SL) and hierarchical FL (HierFL) to reduce device-side computation and improve scalability, but this introduces training complexity due to coordination across tiers. To address these issues, we propose SHeRL-FL, which integrates SL and hierarchical model aggregation and incorporates representation learning at intermediate layers. By allowing clients and edge servers to compute training objectives independently of the cloud, SHeRL-FL significantly reduces both coordination complexity and communication overhead. To evaluate the effectiveness and efficiency of SHeRL-FL, we performed experiments on image classification tasks using CIFAR-10, CIFAR-100, and HAM10000 with AlexNet, ResNet-18, and ResNet-50 in both IID and non-IID settings. In addition, we evaluate performance on image segmentation tasks using the ISIC-2018 dataset with a ResNet-50-based U-Net. Experimental results demonstrate that SHeRL-FL reduces data transmission by over 90\% compared to centralized FL and HierFL, and by 50\% compared to SplitFed, which is a hybrid of FL and SL, and further improves hierarchical split learning methods. 
\end{abstract}

% Uncomment the following to link to your code, datasets, an extended version or similar.
% You must keep this block between (not within) the abstract and the main body of the paper.
% \begin{links}
%     \link{Code}{https://aaai.org/example/code}
%     \link{Datasets}{https://aaai.org/example/datasets}
%     \link{Extended version}{https://aaai.org/example/extended-version}
% \end{links}

\section{Introduction}

    Federated Learning (FL) \cite{konečný2015federatedoptimizationdistributedoptimizationdatacenter, mcmahan2017communication} enables multiple clients to collaboratively train a shared model without exchanging raw data, preserving privacy and data sovereignty in domains such as healthcare, finance, and mobile computing. Despite its advantages, large-scale FL suffers from high communication costs and heavy computation on resource-constrained clients. To alleviate these limitations, hierarchical architectures have been proposed to improve communication efficiency by introducing intermediate edge servers. However, even with hierarchical FL, the computational burden on clients remains a major challenge. To address this, Split Learning (SL) \cite{GUPTA20181} reduces client computation by partitioning a deep neural network between clients and a server. In SL, the client-side model processes local data and sends intermediate activations (smashed data) to the server, which completes the forward pass, computes the loss, and returns gradients for backpropagation. This approach leverages server computation to reduce client workload and has shown promise in various applications \cite{vepakomma2019split, poirot2019split, li2023convergence}. However, SL suffers from sequential bottlenecks, catastrophic forgetting \cite{doi:10.1073/pnas.1611835114} and prolonged client idle time, both of which degrade overall training efficiency.


     To combine the efficiency of SL and the scalability of FL, \cite{thapa2022splitfed} proposed SplitFed, which integrates model partitioning through parallel training. Each client maintains a local copy of the client-side partition, processes private data, and sends smashed data to a central server that hosts the server-side partition. The server completes the forward and backward passes, returning gradients to clients for local updates. SplitFed allows parallel client training, improving scalability while retaining SL’s reduced client computation. However, SplitFed suffers from high client–server communication and synchronization delays \cite{thapa2022splitfed}. Hierarchical Split Federated Learning (HSFL) \cite{10980018} extends this design by adding the edge tier to lower communication costs, relieve device-level constraints, and accelerate convergence. Yet, HSFL often splits models by device capability rather than layer semantics, potentially hindering representation consistency. Independent tier-wise aggregation can destabilize training, and frequent gradient exchanges raise privacy concerns, such as gradient inversion attacks \cite{10025466, guo2025exploringvulnerabilitiesfederatedlearning}.

     \textbf{Motivation.} While recent FL frameworks improve scalability and communication, the quality of intermediate representations exchanged between tiers remains overlooked. Poorly structured features disrupt consistency, slow convergence, and degrade generalization. Representation learning can improve semantic alignment across segmented models, improve information flow, and stabilize aggregation. However, current split-based methods pay limited attention to updating intermediate features, as they are not treated as explicit learning objectives. In this paper, we propose \textbf{SHeRL-FL} (Split and Hierarchical Representation Learning), a role-aware FL architecture that integrates explicit representation learning at edge servers. By aligning model segments with their functional roles, SHeRL-FL enables tier-wise optimization without backward gradients from the cloud, reducing communication and privacy risk.

    \textbf{Contribution.} In summary, the main contributions are as follows:
    \begin{itemize}
        \item We propose SHeRL-FL, a split hierarchical federated learning framework that integrates functional alignment into each tier. Unlike prior approaches that treat tiers uniformly or require end-to-end gradients, SHeRL-FL enables each tier to learn task-relevant representations independently while preserving inter-tier coherence.
        \item By assigning distinct objectives to client, edge, and cloud models, SHeRL-FL improves stability, interpretability, and modularity, making it well-suited for resource-constrained and privacy-sensitive settings.
        \item SHeRL-FL eliminates cloud-to-client gradient transmission, reducing communication costs and mitigating gradient-based privacy risks, which is critical for domains like healthcare and personal analytics.
        \item Experiments on CIFAR-10 ~\cite{krizhevsky2009learning}, CIFAR-100 ~\cite{krizhevsky2009learning}, HAM10000 ~\cite{tschandl2018ham10000}, and ISIC-2018 ~\cite{8363547} across multiple hierarchical settings show consistent gains in accuracy, efficiency, and representation quality over baselines such as FedAvg, FedSGD, FedNova, FedProx, SplitFed, HierFL, and HSFL.
    \end{itemize}


\section{Background and Related Works}

    Conventional FL algorithms include Federated Stochastic Gradient Descent (FedSGD) and Federated Averaging (FedAvg) \cite{mcmahan2017communication}. FedSGD aggregates after each mini-batch, causing frequent communication, while FedAvg improves efficiency by allowing multiple local updates before aggregation, better handling non-IID data. To address heterogeneity, FedProx \cite{MLSYS2020_1f5fe839} adds a proximal term for stable convergence. In addition, FedNova \cite{10.5555/3495724.3496362} tackles client drift in heterogeneous systems by normalizing local updates. 
 
     SL introduced by \cite{GUPTA20181, vepakomma2019split} partitions deep neural networks $\mathbf{W}$ between clients and a server. Clients send smashed data, rather than raw data, to the server, reducing model leakage \cite{9756883}. \cite{10.1145/3320269.3384740} proposed Differential Privacy (DP) incorporated into SL, integrating privacy guarantees while maintaining the utility of the model. However, traditional SL clients follow sequential training processes, resulting in a bottleneck on the server side, particularly on a large scale. SplitFed \cite{thapa2022splitfed} addresses this by allowing clients to train in parallel and send activations to a server, which trains the remaining model layers. Client models are periodically aggregated to maintain synchronized learning. Other variants further improve privacy, scalability, and latency, including SplitNN with partial sharing \cite{vepakomma2019split} and Asynchronous Federated SL \cite{10707222}, and domain-specific applications in healthcare \cite{10746501}, IoT anomaly detection \cite{bdcc8030021}, and cloud-edge recommendation systems \cite{10.1186/s13677-023-00435-5}, demonstrating its privacy-preserving benefits and adaptability across domains.

     To reduce cloud communication at scale, Hierarchical Federated Learning (HierFL) introduces intermediate aggregators (e.g., edge servers)\cite{9148862}. Clients send updates to edge nodes that perform partial aggregation before being sent to the cloud for final aggregation. This hierarchy reduces upstream communication and distributes the server-side computation. Extensions of this work include Hierarchical Global Asynchronous Federated Learning Across Multi-Center \cite{pmlr-v260-xie25c} with straggler-minimizing asynchronous aggregation and Hierarchical Clustering-based Personalized Federated Learning \cite{10.1145/3580795} for hierarchical clustering with non-IID personalization. However, existing HierFL frameworks overlook edge device heterogeneity and rely on client-side full-model training, which burdens constrained devices. This motivates the integration of computation-aware model splitting in a hierarchical structure. Hence, HSFL \cite{10980018} splits the model between clients, edge servers, and the cloud. \cite{10251444} further improves efficiency with distributed training and aggregation. However, most HSFL methods neglect functional tier-specific roles, leading to misaligned representations and suboptimal learning.

     Representation learning offers a potential solution to misalignment by guiding model segments to learn semantically meaningful features. Contrastive learning methods, including SimCLR \cite{10.5555/3524938.3525087} and SupCon \cite{10.5555/3495724.3497291}, have shown success in centralized settings, while ConPro \cite{Nguyen_2024_CVPR} and SemiSE \cite{10980996} propose role-conditioned and semi-supervised strategies for the representation of medical images. Although federated representation learning enhances generalization and personalization under non-IID data, its role in HSFL is still underexplored. Our work bridges this gap by integrating role-aware representation learning into HSFL to support more efficient, modular, and privacy-preserving training.

\section{The Proposed Framework}
    \begin{figure*}[!ht]
        \centering
        \includegraphics[width=\linewidth]{Figures/method.pdf}
        \caption{Overview of FL paradigms. (A) Conventional FL exchanges full models with the cloud. (B) SplitFL transmits smashed data but depends on cloud gradients. (C) Hierarchical FL uses edge servers for intermediate aggregation. (D) Hierarchical SplitFL combines model splitting with hierarchical aggregation. (E) SHeRL-FL (ours) enhances (D) with semantic-aware splitting, edge-level contrastive learning, and gradient-free communication to improve accuracy, efficiency, and scalability.}
        \label{fig:framework}
    \end{figure*}
    We propose \textit{Split and Hierarchical Representation Learning} (SHeRL-FL) to reduce communication and enhance representation alignment in federated systems. As shown in Figure~\ref{fig:framework} and Algorithm~\ref{alg:sherlfl}, SHeRL-FL operates in three tiers: clients, edge servers, and cloud. Training has three stages: (1) clients extract features and send smashed data to edges; (2) edges refine representations via contrastive learning and aggregate locally; (3) the cloud performs final inference and global aggregation.
    
     \textbf{Model Partitioning.} Unlike uniform partitioning in SL, our method aligns the splits with functional roles and system constraints. The client hosts shallow layers for low-level features (edges, textures, shapes), enabling lightweight, privacy-preserving computation. The edge server contains deeper layers for high-level representations and handles representation learning and local aggregation. The cloud holds the task-specific module (classifier/decoder) for final inference and global aggregation, ensuring cross-silo consistency. This function-aware design matches resources to roles, reducing redundancy and improving efficiency.
    
     \textbf{Training Workflow.} Each client processes raw input data using its local model segment and sends the smashed data output to its assigned edge server. Instead of processing these features independently, the edge server forms input pairs and applies contrastive learning to organize the feature space better. We adopt a margin-based contrastive loss that encourages similar samples to have similar representations and dissimilar samples to be far apart in the embedding space. Specifically, the loss is computed as follows:
    \begin{align}
        \mathcal{L}_{\text{contrastive}} &= \mathbb{E}\Big[ y_{ij} \, d_{\mathrm{cos}}(c_i, c_j) \\
        &\quad + (1 - y_{ij}) \, \max\big(0, m - d_{\mathrm{cos}}(c_i, c_j)\big) \Big] \nonumber
    \end{align}
    
    \begin{equation}
        d_{\text{cos}}(c_i, c_j) = \frac{c_i^T c_j}{\|c_i\| \|c_j\|},
    \end{equation}
     \noindent where $y_{ij} = 0$ if $c_i$ and $c_j$ belong to the same class, $y_{ij} = 1$ otherwise, and $m \in [0, 2]$ is a margin hyperparameter. The cosine distance encourages direction-based similarity, which is more suitable for representation learning compared to the Euclidean distance. The contrastive loss is backpropagated by the edges and partially into the clients, enabling immediate local updates without cloud interaction. This decoupling significantly reduces communication overhead and round-trip delays.
    
    \textbf{Hierarchical Model Aggregation.} At intervals $t_1$, each edge aggregates client updates via averaging or similar schemes and then forwards results to the cloud. The cloud computes task loss (e.g., cross-entropy $\mathcal{L}_{\text{task}}$) and performs global aggregation at $t_2$ across edges. This two-level process reduces client–cloud traffic, alleviates communication bottlenecks, enables faster edge updates, and optimizes representations locally before global alignment.

    \textbf{Advantages.} The framework offers key benefits for large-scale FL. Localized training and deferred cloud communication reduce bandwidth and latency. Clients send only intermediate features, improving privacy and reducing leakage risk. Edge-level contrastive learning improves feature space structure and aligns objectives within silos. Hierarchical aggregation balances workload across tiers, scaling to hundreds of clients. Local contrastive learning also boosts robustness under non-IID data by promoting intra-silo alignment.
    \begin{algorithm}[!ht]
        \caption{SHeRL-FL Training}
        \label{alg:sherlfl}
        \begin{algorithmic}[1]
            \State \textbf{Input:} Number of clients $N$, edge servers $E$, cloud server $C$, Number of epochs $E_p$, aggregation intervals $t_1, t_2$, client-side model $M_c$, edge-side model $M_e$, cloud-side model $M_{\mathrm{cloud}}$.
            \State \textbf{Output:} Trained $M_c, M_e, M_{\mathrm{cloud}}$
            \State Initialize $M_c^i, M_e^j, M_{\mathrm{cloud}}$ and optimizers
            \For{$epoch = 1$ to $E_p$}
                \State Select active clients $U \subseteq N$
                \For{client $i \in U$}
                    \State Load $D_i=\{(x^i_1,x^i_2,y^i_1,y^i_2)\}$
                    \State $(f^i_1,f^i_2)\gets M_c^i(x^i_1,x^i_2)$
                    \State Send $(f^i_1,f^i_2,y^i_1,y^i_2)$ to $e_j$
                \EndFor
                \For{edge server $e_j \in E$ }
                    \State Receive smashed data from connected clients  
                    \State Forward it through edge-side mode
                    \State Train $M_e^j$ (contrastive loss)
                    \State Send back gradient to clients to update
                    \State Forward to $M_{\mathrm{cloud}}$
                \EndFor
                \State Cloud trains $M_{\mathrm{cloud}}$ (task loss)
                \If{$epoch \bmod t_1=0$}
                    \State Edge aggregation: update clients
                \EndIf
                \If{$epoch \bmod t_2=0$}
                    \State Cloud aggregation: update edges
                \EndIf
                \State Evaluate $(M_c,M_e,M_{\mathrm{cloud}})$
            \EndFor
            \State \Return Best weights
        \end{algorithmic}
    \end{algorithm}
    
     In general, our framework presented in Algorithm~\ref{alg:sherlfl} combines semantic-aware model splitting, edge-level representation learning, and hierarchical aggregation to deliver a scalable and communication-efficient solution for federated learning in resource-constrained and heterogeneous environments. By reducing communication volume, optimizing resource utilization, and enabling localized processing, SHeRL-FL not only improves performance but also lowers energy consumption and associated carbon emissions, contributing to more environmentally sustainable AI deployment.

\section{Experiments}
    We evaluated on CIFAR-10, CIFAR-100, and HAM10000 for classification (AlexNet, ResNet-18, ResNet-50) and ISIC-2018 for segmentation (ResNet-50 U-Net). Baselines include FedAvg, FedSGD, FedNova, FedProx, HierFL, SplitFL, and HSFL under IID/non-IID. Classification uses cross-entropy loss with the F1 score while segmentation uses IoU and DICE. The models are trained with Adam ($1\mathrm{e}{-4}$ LR) for fixed communication rounds. Traditional FL/SplitFL simulates 200 clients; HierFL, HSFL, and ours use 200 clients and 10 edges. We measure communication cost (data volume) and computation (avg. client RAM). The experiments were run in PyTorch 2.7.0+cu128, Python 3.9.21, on a NVIDIA RTX 3080 Ti.
    \begin{table*}[!ht]
        \centering
        \begin{tabular}{l|cccccccc|c}
            \hline
            \textbf{Method} &
            \begin{tabular}[c]{@{}c@{}}\textbf{Client}\\ \textbf{smashed}\\ \textbf{upload}\end{tabular} &
            \begin{tabular}[c]{@{}c@{}}\textbf{Edge}\\ \textbf{smashed}\\ \textbf{upload}\end{tabular} &
            \begin{tabular}[c]{@{}c@{}}\textbf{Client}\\ \textbf{model}\\ \textbf{upload}\end{tabular} &
            \begin{tabular}[c]{@{}c@{}}\textbf{Client}\\ \textbf{model}\\ \textbf{download}\end{tabular} &
            \begin{tabular}[c]{@{}c@{}}\textbf{Edge}\\ \textbf{model}\\ \textbf{upload}\end{tabular} &
            \begin{tabular}[c]{@{}c@{}}\textbf{Edge}\\ \textbf{model}\\ \textbf{download}\end{tabular} &
            \begin{tabular}[c]{@{}c@{}}\textbf{Gradient}\\ \textbf{from}\\ \textbf{cloud}\end{tabular} &
            \begin{tabular}[c]{@{}c@{}}\textbf{Gradient}\\ \textbf{from}\\ \textbf{edge}\end{tabular} &
            \textbf{Total (GB)} \\
            \hline
            FedAvg & -- & -- & 351.40 & 3,513.95 & -- & -- & -- & -- & 3,865.35 \\
            FedSGD & -- & -- & 351.40 & 3,513.95 & -- & -- & -- & -- & 3,865.35 \\
            FedNova & -- & -- & 351.40 & 3,513.95 & -- & -- & -- & -- & 3,865.35 \\
            FedProx & -- & -- & 351.40 & 3,513.95 & -- & -- & -- & -- & 3,865.35 \\ \hline
            HierFL & -- & -- & 351.40 & 3,865.34 & 175.70 & -- & -- & -- & 4,392.44 \\ \hline
            SplitFed & 6.87 & -- & 1.09 & 216.82 & -- & -- & 3.73 & -- & 228.51 \\ \hline
            HSFL (A) & 54.93 & 13.74 & 0.22 & 0.93 & 1.34 & 1.34 & 12.22 & 31.25 & 115.97 \\
            HSFL (B) & 54.93 & 13.74 & 0.11 & 0.48 & 0.67 & 0.67 & 12.22 & 31.25 & 114.07 \\
            HSFL (C) & 54.93 & 13.74 & 0.06 & 0.24 & 0.34 & 0.34 & 12.22 & 31.25 & 113.12 \\ \hline
            \textbf{Ours (A)} & \textbf{54.93} & \textbf{13.74} & \textbf{0.22} & \textbf{0.93} & \textbf{1.34} & \textbf{1.34} & -- & \textbf{31.25} & \textbf{103.75} \\
            \textbf{Ours (B)} & \textbf{54.93} & \textbf{13.74} & \textbf{0.11} & \textbf{0.48} & \textbf{0.67} & \textbf{0.67} & -- & \textbf{31.25} & \textbf{101.85} \\
            \textbf{Ours (C)} & \textbf{54.93} & \textbf{13.74} & \textbf{0.06} & \textbf{0.24} & \textbf{0.34} & \textbf{0.34} & -- & \textbf{31.25} & \textbf{100.90} \\
            \hline
        \end{tabular}
        \caption{Communication overhead (GB) on CIFAR-100 with ResNet-18 over 200 rounds. 
        (A) is $(t_1=5, t_2=10)$, (B) is $(t_1=10, t_2=20)$, and (C) is $(t_1=25, t_2=50)$.}
        \label{tab:comm_overhead}
    \end{table*}
        \begin{table}
            \centering
            \begin{tabular}{c|c|c}
                \hline
                \textbf{Methods} & \textbf{IoU (\%)} & \textbf{DICE (\%)} \\ \hline
                FedAvg    & 13.73  & 23.38 \\ 
                FedSGD    & 10.23  & 20.54 \\ 
                FedNova   & 19.77  & 25.62 \\ 
                FedProx   & 22.01  & 29.93 \\ \hline
                SplitFed  & 15.26  & 24.46 \\ \hline
                HierFL    & 15.93  & 25.06 \\ \hline
                HSFL (A)  & 24.45 & 32.20 \\
                HSFL (B)   & 24.12 & 32.05 \\
                HSFL (C)     & 23.88 & 31.74 \\ \hline
                Ours (A)  & 28.30 & 35.10 \\
                Ours (B)   & 27.05 & 34.85 \\
                Ours (C)     & 26.72 & 34.42 \\ \hline
            \end{tabular}
            \caption{Segmentation performance comparison of different FL methods in terms of IoU and DICE scores. (A) is $(t_1=5, t_2=10)$, (B) is $(t_1=10, t_2=20)$, and (C) is $(t_1=25, t_2=50)$.}
            \label{tab:sgmentation}
        \end{table}
        \begin{table*}
            \centering
            \begin{tabular}{l|cc|cc|cc}
                \hline
                \multirow{2}{*}{\textbf{Method}} &
                \multicolumn{2}{c|}{\textbf{CIFAR-10}} &
                \multicolumn{2}{c|}{\textbf{CIFAR-100}} &
                \multicolumn{2}{c}{\textbf{HAM10000}} \\
                \cline{2-7}
                 & \textbf{IID} & \textbf{Non-IID} & \textbf{IID} & \textbf{Non-IID} & \textbf{IID} & \textbf{Non-IID} \\
                \hline
                HSFL (baseline) & 49.18\% & 22.58\% & 22.80\% & 0.20\% & 62.28\% & 24.75\% \\
                Ours (A) & 52.64\% & 23.20\% & 25.47\% & 0.53\% & 65.12\% & 25.41\% \\
                Ours (B) & 57.11\% & 37.55\% & 29.73\% & 3.47\% & 69.38\% & 33.58\% \\
                Ours (C) & 60.50\% & 39.62\% & 32.86\% & 4.46\% & 71.57\% & 35.39\% \\
                Ours (D) & 65.44\% & 45.52\% & 55.23\% & 6.65\% & 78.86\% & 38.36\% \\
                \hline
            \end{tabular}
            \caption{Ablation on CIFAR-10, CIFAR-100, and HAM10000 (IID/Non-IID) with AlexNet, ResNet-18, ResNet-50. HSFL baseline vs SHeRL-FL components in F1 Score at $t_1=5$, $t_2=10$. (A) SHeRL-FL without contrastive loss and role-aware splitting, (B) with contrastive loss, (C) with role-aware splitting, and (D) full configuration.}
            \label{tab:ablation}
        \end{table*}
        \begin{table}[!ht]
            \centering
            \begin{tabular}{l|c|c}
            \hline
            \centering\textbf{Methods} & \textbf{IoU (\%)} & \textbf{DICE (\%)} \\ \hline
            HSFL (baseline) & 24.45 & 32.20 \\
            Ours (A) & 25.18 & 33.04 \\
            Ours (B) & 26.02 & 34.12 \\
            Ours (C) & 26.63 & 34.79 \\
            Ours (D)& 28.30 & 35.10 \\
            \hline
            \end{tabular}
            \caption{Ablation study on ISIC-2018 (IID) with ResNet-50 U-Net, comparing the HSFL baseline and SHeRL-FL variants in IoU and DICE at $t_1=5$ and $t_2=10$. (A) SHeRL-FL without contrastive loss and role-aware splitting, (B) with contrastive loss, (C) with role-aware splitting, and (D) full configuration.}
            \label{tab:ablation_segmentation}
        \end{table}
        \begin{table*}[!ht]
            \centering
            \begin{tabular}{l|cc|cc|cc}
                \hline
                \multirow{2}{*}{\textbf{$m$}} &
                \multicolumn{2}{c|}{\textbf{CIFAR-10}} &
                \multicolumn{2}{c|}{\textbf{CIFAR-100}} &
                \multicolumn{2}{c}{\textbf{HAM10000}} \\ \cline{2-7}
                 & \textbf{IID} & \textbf{Non-IID} & \textbf{IID} & \textbf{Non-IID} & \textbf{IID} & \textbf{Non-IID} \\
                \hline
                0.2 & 54.28\% & 36.00\% & 48.15\% & 25.10\% & 71.42\% & 30.55\% \\
                0.5 & 65.44\% & 45.52\% & 55.23\% & 36.65\% & 78.86\% & 38.36\% \\
                1.0 & 59.77\% & 37.49\% & 52.88\% & 30.12\% & 76.40\% & 34.70\% \\
                1.5 & 60.03\% & 41.99\% & 53.10\% & 32.50\% & 76.85\% & 36.10\% \\
                2.0 & 54.38\% & 40.44\% & 48.65\% & 28.00\% & 72.12\% & 32.45\% \\
                \hline
            \end{tabular}
            \caption{Ablation on CIFAR-10, CIFAR-100, and HAM10000 (IID/Non-IID) with AlexNet, ResNet-18, and ResNet-50. Comparison of different margins $m$ in contrastive loss at edge servers of SHeRL-FL in terms of F1 Score at $t_1=5$, $t_2=10$.}            \label{tab:margin_ablation}
        \end{table*}
        \begin{table}[ht]
            \centering
            \begin{tabular}{l|c|c}
                \hline
                \centering\textbf{$m$} & \textbf{IoU (\%)} & \textbf{DICE (\%)} \\ \hline
                0.2 & 24.50 & 30.80 \\
                0.5 & 28.30 & 35.10 \\
                1.0 & 27.20 & 34.00 \\
                1.5 & 26.80 & 33.60 \\
                2.0 & 25.10 & 31.50 \\
                \hline
            \end{tabular}
            \caption{Ablation on ISIC-2018 (IID) with ResNet-50 U-Net. Comparison of margins $m$ in contrastive loss at edge servers of SHeRL-FL in IoU and DICE at $t_1=5$, $t_2=10$.}
            \label{tab:ablation_segmentation_margin}
        \end{table}
        
     \textbf{F1 Score over training rounds.}
        \begin{figure}[!ht]
            \centering
            \begin{subfigure}[b]{\linewidth}
                \centering
                \includegraphics[width=\linewidth]{Figures/cifar_train_accuracy_comparison.pdf}
                \caption{Under IID setting}
                \label{fig:f1_score_per_round_cifar10_iid}
            \end{subfigure}
            \begin{subfigure}[b]{\linewidth}
                \centering
                \includegraphics[width=\linewidth]{Figures/cifar_train_accuracy_comparison_noniid.pdf}
                \caption{Under non-IID setting}
                \label{fig:f1_score_per_round_cifar10_noniid}
            \end{subfigure}
            \caption{The F1 Score Comparison across methods with the CIFAR-10 dataset and the AlexNet backbone under different settings.}
            \label{fig:f1_score_per_round_cifar10}
        \end{figure}
        \begin{figure}[!ht]
            \centering
            \begin{subfigure}[b]{0.49\linewidth}
                \includegraphics[width=\linewidth]{Figures/tsne_encodediid.pdf}
                \caption{Under IID setting} % <-- caption text optional
                \label{fig:tsne_encodediid}
            \end{subfigure}
            \hfill
            \begin{subfigure}[b]{0.49\linewidth}
                \includegraphics[width=\linewidth]{Figures/tsne_encodednoniid.pdf}
                \caption{Under non-IID setting} % <-- caption text optional
                \label{fig:tsne_encodednoniid}
            \end{subfigure}
            \caption{t-SNE visualizations of learned representations of the CIFAR-10 dataset under different settings.}
            \label{fig:tsne_combined}
        \end{figure}
        \begin{figure}[!ht]
            \centering
            \begin{subfigure}[b]{\linewidth}
                \centering
                \includegraphics[width=\linewidth]{Figures/cifar100_train_accuracy_comparison.pdf}
                \caption{Under IID setting}
                \label{fig:f1_score_per_round_cifar100_iid}
            \end{subfigure}
            \begin{subfigure}[b]{\linewidth}
                \centering
                \includegraphics[width=\linewidth]{Figures/cifar100_train_accuracy_comparison_noniid.pdf}
                \caption{Under non-IID setting}
                \label{fig:f1_score_per_round_cifar100_noniid}
            \end{subfigure}
            \caption{The F1 Score Comparison across methods with the CIFAR-100 dataset and the ResNet-18 backbone under different settings.}
            \label{fig:f1_score_per_round_cifar100} % Move this here, after \caption
        \end{figure}
        \begin{figure}[!ht]
            \centering
            \begin{subfigure}[b]{\linewidth}
                \centering
                \includegraphics[width=\linewidth]{Figures/ham10000_train_accuracy_comparison.pdf}
                \caption{Under IID setting}
                \label{fig:f1_score_per_round_ham10000_iid}
            \end{subfigure}
            \begin{subfigure}[b]{\linewidth}
                \centering
                \includegraphics[width=\linewidth]{Figures/ham10000_train_accuracy_comparison_noniid.pdf}
                \caption{Under non-IID setting}
                \label{fig:f1_score_per_round_ham10000_noniid}
            \end{subfigure}
            \caption{The F1 Score Comparison across methods with the HAM10000 dataset and the ResNet-50 backbone under different settings.}
            \label{fig:f1_score_per_round_ham10000}
        \end{figure}
        In the case of the CIFAR-10 dataset under the IID setting with the AlexNet backbone, Figure~\ref{fig:f1_score_per_round_cifar10}(\subref{fig:f1_score_per_round_cifar10_iid}) shows that our proposed method significantly outperforms all baselines in terms of the F1 score in 200 communication rounds. Specifically, the configuration with synchronization intervals ($t_1 = 10, t_2 = 20$) achieves the highest F1 score, exceeding 65\%. In contrast, traditional methods such as FedAvg, FedSGD, FedNova, and FedProx remain below 55\%, while SplitFed and hierarchical baselines such as HierFL and HSFL converge more slowly and plateau around 60\%. Our method maintains competitive performance even with larger aggregation intervals (e.g., $t_1 = 25, t_2 = 50$). 
        
        In terms of non-IID settings, Figure \ref{fig:f1_score_per_round_cifar10}(\subref{fig:f1_score_per_round_cifar10_noniid}) illustrates our proposed method consistently outperforms all baselines, with the configuration $(t_1=10, t_2=20)$ achieving the highest F1 score. In addition, traditional methods such as FedAvg, FedProx, and FedSGD perform poorly, with slow or stagnant improvements, while SplitFed shows fast initial gains but suffers from instability.
        
        To explain these performances, Figure \ref{fig:tsne_combined} shows the cluster of smashed data from the test CIFAR-10 dataset after going through the edge-side model in IID and non-IID settings. The label follows the distribution of all samples, so the performance of the specific task module in the cloud remains.
        This result highlights the effectiveness and stability of our proposed framework, particularly in scenarios with balanced data distributions and lightweight model architectures.
        
        In both IID (Figure~\ref{fig:f1_score_per_round_cifar100}(\subref{fig:f1_score_per_round_cifar100_iid})) and Non-IID (Figure~\ref{fig:f1_score_per_round_cifar100}(\subref{fig:f1_score_per_round_cifar100_noniid})) settings on CIFAR-100 with ResNet-18, SHeRL-FL consistently outperforms all baselines in F1 score. While FedAvg and FedProx perform reasonably in IID, they degrade sharply in Non-IID. FedNova and FedSGD improve steadily but plateau at lower scores. SplitFed shows instability with large fluctuations, especially in Non-IID. HierFL and all HSPL variants remain near zero in both cases. SHeRL-FL achieves the highest and most stable F1 across rounds, with the best result at $(t_1=10, t_2=20)$, demonstrating robustness to data heterogeneity.
        
        As shown in Figure \ref{fig:f1_score_per_round_ham10000}(\subref{fig:f1_score_per_round_ham10000_iid}) displays the F1 Score in the HAM10000 dataset with IID setting. Specifically, FedAvg and FedProx converge to F1 scores around 70\%. In contrast, SplitFed and HierFL exhibit slower convergence and slightly lower final performance. In addition, HSFL variants show improved results as synchronization intervals ($t_1$, $t_2$) decrease. In particular, the proposed method outperforms all baselines in all configurations, achieving the highest F1 score of approximately 78\% with ($t_1=5$, $ t_2=10$).

        Figure \ref{fig:f1_score_per_round_ham10000}(\subref{fig:f1_score_per_round_ham10000_iid}) presents the results under non-IID conditions. In this setting, the performance of all methods declines relative to the IID scenario. FedSGD exhibits limited improvement over communication rounds, while FedAvg, FedNova, and FedProx show slower and less stable convergence. Furthermore, although HSFL variants benefit from increased synchronization intervals, their performance remains sensitive to data heterogeneity. In comparison, the proposed method consistently achieves the highest F1 scores, with the configuration ($t_1=10$, $t_2=20$) providing the most favorable balance between convergence speed and accuracy.
        
        
     \textbf{Performance of Segmentation task.} Table~\ref{tab:sgmentation} presents the segmentation performance of various baselines of FL and our proposed method, evaluated using IoU and DICE scores. Among traditional FL methods, FedProx achieves the highest accuracy, while SplitFed and HierFL offer marginal improvements due to model splitting and hierarchical model aggregation. HSFL demonstrates significant gains over these baselines, while our method consistently outperforms all compared approaches in all aggregation settings, achieving the best results with $t_1{=}25$ and $t_2{=}50$, highlighting the effectiveness of joint split hierarchy representation learning under low frequency synchronization.
        
     \textbf{Computational cost.} Figure~\ref{fig:ram_per_round} compares average client RAM usage across methods. Our approach maintains stable consumption around 380 MB, well below HierFL ($\sim$650 MB) and flat FL methods such as FedAvg, FedSGD, FedProx, and FedNova ($\sim$580 MB). While SplitFed and HSFL use less memory ($<$150 MB), their accuracy is notably lower. The slight increase over HSFL stems from processing image pairs instead of single images. Overall, our method offers a favorable trade-off, combining strong performance with low, stable memory usage, which is suitable for memory-constrained clients.
        \begin{figure}[!t]
            \centering
            \includegraphics[width=\linewidth]{Figures/cifar_client_gpu_ram_list_comparison.pdf}
            \caption{The RAM Usage at client tier across mothods with the CIFAR-10 dataset and the AlexNet backbone under IID Scenario.}
            \label{fig:ram_per_round}
        \end{figure}
         
        
     \textbf{Communication cost.} Table~\ref{tab:comm_overhead} reports the communication overhead (GB) for 200 clients. Hierarchical methods assume 200 clients linked through 10 edge servers, while flat baselines communicate directly with the cloud. All methods use a 10\% client sampling rate. Flat methods (FedAvg, FedSGD, FedNova, FedProx) incur very high costs ($>$3{,}865 GB) due to full model exchange each round. HierFL reaches the highest total (4{,}392.44 GB) but shifts much of the load from the cloud to edges, benefiting scenarios with limited cloud bandwidth. SplitFed lowers cost (216.82 GB) by uploading smashed data, but still needs gradient feedback from the cloud. HSFL reduces the total to $\sim$115 GB by limiting gradient exchanges, though cloud-to-edge gradients remain. Our method achieves the lowest overhead (100.90 GB) by removing cloud-to-edge gradient transfers; edges coordinate client updates locally. While smashed activations from clients (54.93 GB) and edges (13.74 GB) add volume, these are transmitted incrementally, minimizing peak bandwidth. Overall, the approach offers strong scalability, bandwidth efficiency, and practicality for large-scale deployments with constrained cloud capacity.
        
     \textbf{Ablation Study.} The ablation results (Tables~\ref{tab:ablation}--\ref{tab:ablation_segmentation}) show that both contrastive loss and role-aware splitting are key to SHeRL-FL’s effectiveness. Removing both components offers only marginal gains over HSFL, while adding either component individually provides noticeable improvements, especially under non-IID settings. The best performance is consistently achieved when both are combined, confirming their complementary roles in improving representation alignment, optimizing partitioning, and improving generalization across tasks.

     The margin study (Tables~\ref{tab:margin_ablation}--\ref{tab:ablation_segmentation_margin}) further reveals that a moderate value of $m=0.5$ produces the most stable and accurate results for classification and segmentation. Smaller margins lead to insufficient interclass separation, whereas larger margins can overly penalize feature alignment, reducing model stability. The consistent peak at $m=0.5$ highlights the importance of carefully tuning the contrast loss margin to balance the strength of separation and the stability of optimization.

\section{Comparison with HSFL}
    Compared to HSFL, SHeRL-FL introduces key architectural and training enhancements that lead to both accuracy and efficiency gains. The integration of contrastive learning at the edge aligns feature spaces without requiring end-to-end gradient flow, reducing representation drift and improving performance, particularly in non-IID settings. Role-aware splitting assigns model segments to tiers more systematically, stabilizing training and accelerating convergence. Communication efficiency is improved by eliminating cloud-to-edge gradient transmission, while maintaining comparable or lower bandwidth usage. Although HSFL achieves lower peak memory usage by processing single images, SHeRL-FL balances slightly higher memory consumption with consistently superior accuracy, faster convergence, and more flexible synchronization. These advantages make SHeRL-FL a more practical choice for large-scale, heterogeneous federated deployments.

\section{Limitations}
    The current role‑aware split is fixed and heuristic; adaptive partitioning could improve efficiency in varied environments. Contrastive learning assumes labeled edge data, which may limit applicability; self‑supervised alternatives are a potential extension. While communication volume is reduced, real network latency tests would further strengthen deployment practicality.
\section{Conclusion}
SHeRL-FL enhances representation learning flexibility while delivering practical communication and memory gains in federated settings. By embedding contrastive learning at the edge, it reduces reliance on cloud synchronization, enabling asynchronous updates and lowering idle time, particularly in heterogeneous or bandwidth-limited environments. Although the communication reduction over HSFL is modest, SHeRL-FL supports more flexible scheduling under real-world constraints. Experiments show consistent accuracy improvements in both IID and non-IID settings, with the largest gains in non-IID cases due to contrastive loss aligning feature spaces. Role-aware splitting further stabilizes tier training and mitigates representation drift, leading to faster convergence. The removal of cloud-to-edge gradient transfers improves bandwidth efficiency and allows more adaptive synchronization, making SHeRL-FL well suited for large-scale, resource-constrained deployments.




\bibliography{aaai2026}

% Check whether the conference requires a reproducibility checklist to be included in the paper.
% If so, you can uncomment the following line and ajust the path to include it.
\makeatletter
\@ifundefined{isChecklistMainFile}{
  % We are compiling a standalone document
  \newif\ifreproStandalone
  \reproStandalonetrue
}{
  % We are being \input into the main paper
  \newif\ifreproStandalone
  \reproStandalonefalse
}
\makeatother

\ifreproStandalone
\documentclass[letterpaper]{article}
\usepackage[submission]{aaai2026}
\setlength{\pdfpagewidth}{8.5in}
\setlength{\pdfpageheight}{11in}
\usepackage{times}
\usepackage{helvet}
\usepackage{courier}
\usepackage{xcolor}
\frenchspacing

\begin{document}
\fi
\setlength{\leftmargini}{20pt}
\makeatletter\def\@listi{\leftmargin\leftmargini \topsep .5em \parsep .5em \itemsep .5em}
\def\@listii{\leftmargin\leftmarginii \labelwidth\leftmarginii \advance\labelwidth-\labelsep \topsep .4em \parsep .4em \itemsep .4em}
\def\@listiii{\leftmargin\leftmarginiii \labelwidth\leftmarginiii \advance\labelwidth-\labelsep \topsep .4em \parsep .4em \itemsep .4em}\makeatother

\setcounter{secnumdepth}{0}
\renewcommand\thesubsection{\arabic{subsection}}
\renewcommand\labelenumi{\thesubsection.\arabic{enumi}}

\newcounter{checksubsection}
\newcounter{checkitem}[checksubsection]

\newcommand{\checksubsection}[1]{%
  \refstepcounter{checksubsection}%
  \paragraph{\arabic{checksubsection}. #1}%
  \setcounter{checkitem}{0}%
}

\newcommand{\checkitem}{%
  \refstepcounter{checkitem}%
  \item[\arabic{checksubsection}.\arabic{checkitem}.]%
}
\newcommand{\question}[2]{\normalcolor\checkitem #1 #2 \color{blue}}
\newcommand{\ifyespoints}[1]{\makebox[0pt][l]{\hspace{-15pt}\normalcolor #1}}

\section*{Reproducibility Checklist}

\vspace{1em}
\hrule
\vspace{1em}

\textbf{Instructions for Authors:}

This document outlines key aspects for assessing reproducibility. Please provide your input by editing this \texttt{.tex} file directly.

For each question (that applies), replace the ``Type your response here'' text with your answer.

\vspace{1em}
\noindent
\textbf{Example:} If a question appears as
%
\begin{center}
\noindent
\begin{minipage}{.9\linewidth}
\ttfamily\raggedright
\string\question \{Proofs of all novel claims are included\} \{(yes/partial/no)\} \\
Type your response here
\end{minipage}
\end{center}
you would change it to:
\begin{center}
\noindent
\begin{minipage}{.9\linewidth}
\ttfamily\raggedright
\string\question \{Proofs of all novel claims are included\} \{(yes/partial/no)\} \\
yes
\end{minipage}
\end{center}
%
Please make sure to:
\begin{itemize}\setlength{\itemsep}{.1em}
\item Replace ONLY the ``Type your response here'' text and nothing else.
\item Use one of the options listed for that question (e.g., \textbf{yes}, \textbf{no}, \textbf{partial}, or \textbf{NA}).
\item \textbf{Not} modify any other part of the \texttt{\string\question} command or any other lines in this document.\\
\end{itemize}

You can \texttt{\string\input} this .tex file right before \texttt{\string\end\{document\}} of your main file or compile it as a stand-alone document. Check the instructions on your conference's website to see if you will be asked to provide this checklist with your paper or separately.

\vspace{1em}
\hrule
\vspace{1em}

% The questions start here

\checksubsection{General Paper Structure}
\begin{itemize}

\question{Includes a conceptual outline and/or pseudocode description of AI methods introduced}{(yes/partial/no/NA)}
yes
\question{Clearly delineates statements that are opinions, hypothesis, and speculation from objective facts and results}{(yes/no)}
yes

\question{Provides well-marked pedagogical references for less-familiar readers to gain background necessary to replicate the paper}{(yes/no)}
yes

\end{itemize}
\checksubsection{Theoretical Contributions}
\begin{itemize}

\question{Does this paper make theoretical contributions?}{(yes/no)}
no

	\ifyespoints{\vspace{1.2em}If yes, please address the following points:}
        \begin{itemize}
	
	\question{All assumptions and restrictions are stated clearly and formally}{(yes/partial/no)}
	Type your response here

	\question{All novel claims are stated formally (e.g., in theorem statements)}{(yes/partial/no)}
	Type your response here

	\question{Proofs of all novel claims are included}{(yes/partial/no)}
	Type your response here

	\question{Proof sketches or intuitions are given for complex and/or novel results}{(yes/partial/no)}
	Type your response here

	\question{Appropriate citations to theoretical tools used are given}{(yes/partial/no)}
	Type your response here

	\question{All theoretical claims are demonstrated empirically to hold}{(yes/partial/no/NA)}
	Type your response here

	\question{All experimental code used to eliminate or disprove claims is included}{(yes/no/NA)}
	Type your response here
	
	\end{itemize}
\end{itemize}

\checksubsection{Dataset Usage}
\begin{itemize}

\question{Does this paper rely on one or more datasets?}{(yes/no)}
yes

\ifyespoints{If yes, please address the following points:}
\begin{itemize}

	\question{A motivation is given for why the experiments are conducted on the selected datasets}{(yes/partial/no/NA)}
	yes

	\question{All novel datasets introduced in this paper are included in a data appendix}{(yes/partial/no/NA)}
	NA

	\question{All novel datasets introduced in this paper will be made publicly available upon publication of the paper with a license that allows free usage for research purposes}{(yes/partial/no/NA)}
	NA

	\question{All datasets drawn from the existing literature (potentially including authors' own previously published work) are accompanied by appropriate citations}{(yes/no/NA)}
	yes

	\question{All datasets drawn from the existing literature (potentially including authors' own previously published work) are publicly available}{(yes/partial/no/NA)}
	yes

	\question{All datasets that are not publicly available are described in detail, with explanation why publicly available alternatives are not scientifically satisficing}{(yes/partial/no/NA)}
	NA

\end{itemize}
\end{itemize}

\checksubsection{Computational Experiments}
\begin{itemize}

\question{Does this paper include computational experiments?}{(yes/no)}
yes

\ifyespoints{If yes, please address the following points:}
\begin{itemize}

	\question{This paper states the number and range of values tried per (hyper-) parameter during development of the paper, along with the criterion used for selecting the final parameter setting}{(yes/partial/no/NA)}
	yes

	\question{Any code required for pre-processing data is included in the appendix}{(yes/partial/no)}
	no

	\question{All source code required for conducting and analyzing the experiments is included in a code appendix}{(yes/partial/no)}
	no

	\question{All source code required for conducting and analyzing the experiments will be made publicly available upon publication of the paper with a license that allows free usage for research purposes}{(yes/partial/no)}
	yes
        
	\question{All source code implementing new methods have comments detailing the implementation, with references to the paper where each step comes from}{(yes/partial/no)}
	yes

	\question{If an algorithm depends on randomness, then the method used for setting seeds is described in a way sufficient to allow replication of results}{(yes/partial/no/NA)}
	yes

	\question{This paper specifies the computing infrastructure used for running experiments (hardware and software), including GPU/CPU models; amount of memory; operating system; names and versions of relevant software libraries and frameworks}{(yes/partial/no)}
	yes

	\question{This paper formally describes evaluation metrics used and explains the motivation for choosing these metrics}{(yes/partial/no)}
	yes

	\question{This paper states the number of algorithm runs used to compute each reported result}{(yes/no)}
	yes

	\question{Analysis of experiments goes beyond single-dimensional summaries of performance (e.g., average; median) to include measures of variation, confidence, or other distributional information}{(yes/no)}
	no

	\question{The significance of any improvement or decrease in performance is judged using appropriate statistical tests (e.g., Wilcoxon signed-rank)}{(yes/partial/no)}
	yes

	\question{This paper lists all final (hyper-)parameters used for each model/algorithm in the paper’s experiments}{(yes/partial/no/NA)}
	yes


\end{itemize}
\end{itemize}
\ifreproStandalone
\end{document}
\fi

\end{document}
